\chapter{Gynecomastia (Enlarged Male Breast Tissue)}

\section*{Definition}
Gynecomastia is benign enlargement of male breast tissue caused by an imbalance of estrogen to androgen activity. It affects up to 65 percent of adolescent boys and 30-50 percent of older men at some point in their lives.

\section*{What Happens in Your Body}
Normal male breast tissue is maintained by a balance between androgens and estrogens. Gynecomastia develops when there is an increased estrogen-to-androgen ratio, either from increased estrogen production, decreased androgen production, or increased peripheral conversion of androgens to estrogens by aromatase enzyme. Histologically, there is ductal and stromal hyperplasia.

\section*{Causes}
\begin{itemize}
\item Physiologic (neonatal, pubertal, aging)
\item Medication-induced (spironolactone, cimetidine, calcium channel blockers, anabolic steroids, alcohol)
\item Liver cirrhosis (impaired estrogen metabolism)
\item Testosterone deficiency (primary or secondary hypogonadism)
\item Testicular tumors (Leydig cell, Sertoli cell, germ cell)
\item Adrenal tumors (estrogen-secreting)
\item Hyperthyroidism
\item Obesity (increased aromatase activity in fat)
\item Chronic kidney disease or dialysis
\item Malnutrition or refeeding
\end{itemize}

\section*{When to See a Doctor}
\begin{itemize}
\item Enlargement of one or both breasts in a male
\item Palpable glandular breast tissue (firm, rubbery disc) centered under the nipple
\item Tenderness or sensitivity of the breast tissue
\item May be unilateral or bilateral
\item No skin changes, retraction, or nipple discharge
\item Slow or gradual growth
\end{itemize}

\section*{Related Symptoms}
\begin{itemize}
\item Breast tenderness in males
\item Testicular pain or swelling
\item Hypogonadism (low testosterone)
\item Breast cancer in males (rare)
\item Pseudogynecomastia (enlargement from fat, not glandular tissue)
\end{itemize}

\section*{Self-Care/Home Management}
\begin{itemize}
\item Reassurance for mild, physiologic cases
\item Lose weight if overweight (reduces aromatase activity)
\item Avoid anabolic steroids, alcohol, marijuana
\item Discontinue offending medications if possible
\item Wear loose-fitting clothing to reduce discomfort
\end{itemize}

\section*{How Your Doctor Will Evaluate You}
\begin{itemize}
\item History and physical examination (glandular vs. fatty tissue)
\item CBC, liver function, renal function, TSH
\item Testosterone, estradiol, LH, FSH levels
\item hCG and tumor markers if testicular tumor suspected
\item Mammography or ultrasound to distinguish from breast cancer
\item Testicular ultrasound if exam is abnormal
\end{itemize}

\section*{Treatment Options}
\begin{itemize}
\item Discontinue or change offending medication
\item Treat underlying cause (hypogonadism, hyperthyroidism)
\item Selective estrogen receptor modulators (tamoxifen)
\item Aromatase inhibitors (anastrozole)
\item Surgery (subcutaneous mastectomy, liposuction) for persistent gynecomastia
\item Observation for pubertal gynecomastia (spontaneous resolution in 90 percent)
\end{itemize}
